\section{Introduction}
Modeling and forecasting asset return volatility is a cornerstone of quantitative finance, with direct applications in option pricing, risk management, and portfolio optimization. Since the pioneering work of \citet{engle1982autoregressive} and \citet{bollerslev1986generalized}, Autoregressive Conditional Heteroskedasticity (ARCH) and Generalized ARCH (GARCH) models have served as the standard framework for capturing conditional heteroskedasticity and volatility clustering. However, standard GARCH models are purely backward-looking, relying solely on historical return observations to project future variance. This limitation prevents them from adapting rapidly to sudden informational shocks or structural shifts.

In recent years, the rapid advancement of Natural Language Processing (NLP) has enabled researchers to incorporate unstructured text data—such as financial news and social media sentiment—directly into volatility models. Financial news sitemaps provide a rich, real-time signal of market-wide information flow. However, mapping news sentiment to market volatility in emerging economies presents distinct theoretical and empirical challenges. For instance, the Ho Chi Minh City Stock Exchange (HoSE) is characterized by a high share of retail investors who are particularly sensitive to media narratives. Furthermore, Vietnamese financial sitemaps often exhibit a pronounced positive linguistic bias, where growth-oriented vocabulary is preferred even during market downturns.

This paper asks whether a hybrid GARCH + LSTM residual-correction model built from VN-Index prices and CafeF sentiment scores can improve volatility forecasts relative to a baseline GARCH(1,1). Specifically, we implement a two-stage workflow on the Vietnamese stock market index (VN-Index) over the ten-year period from 2015 to 2024. In the first stage, a baseline GARCH(1,1) model captures linear time-series volatility clustering. In the second stage, a Long Short-Term Memory (LSTM) recurrent neural network is trained on GARCH residuals using lagged returns, volatility features, daily article counts, and sentiment aggregates derived from article-level PhoBERT inference.

Our empirical results show that the hybrid model outperforms the baseline GARCH(1,1) model on the held-out test window. The reported experiment reduces RMSE by $14.48\%$ and MAPE by $48.00\%$, and the Diebold-Mariano test rejects equal predictive accuracy at the $99.9\%$ confidence level ($p < 0.001$). The robustness checks suggest that the residual-correction architecture drives most of the gain, while a linear GARCH-X specification does not generalize well.

The remainder of this paper is organized as follows. Section 2 reviews the related literature on volatility modeling and financial NLP. Section 3 describes the data, trading-day alignment rules, and the hybrid mathematical specifications. Section 4 presents the empirical results and diagnostics. Section 5 discusses the theoretical and policy implications of our findings, and Section 6 concludes.
